\chapter{模型边界、未解决问题与可检验的研究方向}
\section{先区分方程缺失与尚未算够}
“没有解决”至少有三种含义：方程根本没有这个自由度；方程包含相关项，但闭合没有经过验证；原则上能够研究，但可用计算的尺度、维数和统计时间不足。这三种情况需要不同的下一步。加密网格不能恢复被方程删去的非热分布；更复杂闭合也不能替代真正的三维尺寸扫描。

\begin{longtable}{P{35mm}P{57mm}P{56mm}}
\toprule 问题 & 当前模型的限制 & 可以如何继续\\\midrule
电子非热尾、束流、速度空间多峰 & 标量／有限矩没有唯一分布函数；相同低阶矩可对应不同尾部 & 动理学对照或明确的重构假设；不从温度直接推断加速谱\\
电子扩散区非对角压力 & 普通Hall MHD缺失；十矩大小仍由热流闭合决定 & 同初值比较局地、梯度、非局地闭合，并做欧姆定律残差\\
粒子熵产生与不可逆热化 & 理想流体能量转移不等于微观熵增；等温模型有隐含热库 & 分物种能量预算；碰撞、相混合和数值耗散分别记账\\
全系统触发及等待时间 & 大种子GEM绕过线性触发；边界可能持续供能或阻碍流出 & 噪声、驱动幅度、片厚和边界的独立扫描\\
电子层长度选择 & 局地标度常仍含未决定的层长和上游场 & 同时测长度、局地场、闭合／耗散及全局率\\
三维持续湍流与X线传播 & 二维排除了某些不稳定性；有限第三维限制可用波数 & 第三维长度与分辨率分别收敛；记录前沿和寿命统计\\\bottomrule
\end{longtable}

\section{闭合能否跨算例使用，是目前最具体的难题}
局地松弛可在一个算例调到接近PIC，却在另一系统产生错误的磁岛、各向异性或增长时间。岛合并的证据已显示：改变离子闭合可以改善厚度与压力支撑，但未必同时修复平均率的尺寸标度。\citep{Ng2015,Ng2017,Ng2020}

2026年的三维结果则更明确：恢复低混合漂移和弯曲模，与扩大到不合理区域的各向异性同时发生。\citep{Closure2026}因此有价值的目标不是“再选一个使速率合适的 $k_0$”，而是提前固定闭合，在Harris片、岛合并、不对称片、导引场与三维扰动中共同检验。应把闭合的参数训练算例与验证算例分开。

本报告建议的可判别指标是：增长率与模数、电子／离子压力张量误差、X点欧姆定律分量、混合宽度、峰值和时间平均重联率，以及闭合所产生的能量通量。只有速率接近而压力错误，不能称为模型等价。

\section{“普适0.1”究竟是哪一个量的普适性}
GEM、非对称Hall和岛合并使用的初值、上游定义与统计窗不同。无显式耗散惯性模型的有限Hall扫描也得到速率随参数变化。\citep{Birn2001,Cassak2008,Ng2015,Andres2016}有待澄清的不是一个孤立小数，而是：哪些局地几何或压力约束能选择开口角；这些约束是否对外部驱动、系统尺寸与闭合保持稳定。

可行的Hall／双流体研究应同时给出有量纲电场、全局初始参数归一化和局地上游归一化；先固定 $d_i,d_e$ 与闭合，再改变盒子尺寸，避免改变 $L/d_i$ 时也无意改变分辨率或耗散尺度。还应分别检验层长是输出，还是理论中预设的参数。

\section{快慢分支、磁岛循环和三维湍流是否属于同一相图}
小系统中的快慢双稳态、大系统二维中的循环转换、三维中的持续湍流，目前不能仅凭名称合并为一个完整相图。\citep{Cassak2005,Cassak2010,Huang2011,Huang2024}它们的电阻、超扩散、电子惯性、边界、压力闭合和扫描路径不同。

建议选取统一方程，至少以 $S,L/d_i,B_g/B_{\rm rec}$ 为独立轴，并对每一点从快、慢两种历史启动。先在二维定位寿命与转换，再用相同面内网格加入第三维；观察长期单X线、多磁岛循环和持续湍流是否存在重叠参数区。这是两种流体方法原则上能够推进的问题，但对真实电子层物理的推断仍需动理学对照。

\section{强导引场的低频约化何时失效}
强导引场并不允许无限制地把一切简化成两场。约化模型需满足频率、各向异性、密度扰动和磁化尺度排序；有限电子温度时，压力张量同阶修正也可能不可省。Schep等的原文直接指出，相关传播尺度上的Landau共振会导致流体模型失效；后来的细丝化研究也明确承认无限变薄后的失效。\citep{Schep1994,L049}

因此对观察到的细丝、爆发或饱和，应额外记录最小尺度相对于 $d_e,r_{Le}$ 及沿场自由程的位置，比较流体与动理学的能量分配。单靠总能量守恒和精细图像不足以判定近似仍成立。

\section{实验比较应包含仪器和装置几何}
实验中的导引场堆积、流出受限、驱动电路和有限探针空间分辨率，都会改变同一个局地层的可观测信号。MRX导引场实验已显示局地Hall作用与全局边界作用难以分离。\citep{GuideExp2012}有意义的下一步是在模拟中生成与实验同位置、同时间平均、同空间卷积的合成诊断，而不是把未滤波的二维切片与实验单条剖面直接对照。

\section{三维前沿与湍动的启动竞争}
Hall与双流体已经能研究重联传播，但未统一预测非均匀电流片、有限自由磁能和斜向模相互作用下的前沿。可检验的问题是：在保持局地片厚与扰动不变时，独立改变可用磁通、边界和长宽比，X线伸长转成仅平移的条件是否仍由磁能耗尽决定？这需要能量收支和多组参数扫描，而非单幅传播图。\citep{L275}

十矩三维研究还显示，电子密度噪声从$10^{-10}$改到$10^{-2}$，LHDI启动时间可从约$30$变到$4\Omega_i^{-1}$。磁种子驱动的重联与噪声驱动的不稳定性谁先形成，能改变后续湍动与多X线结构。下一步须同时扫描噪声谱、幅度、闭合与网格；闭合在固定噪声下“吻合一次”，并不足以证明对自然启动的预测能力。\citep{E552}

\section{把来源信息引入流体：新模型应怎样接受检验}
两个电子群的研究把“寻找一个普遍热流公式”改成“增加有限数量的可演化来源标签”。已有低$\beta$固定离子证据说明它可能有效，但真正的研究问题是：何时分群、何时合并、怎样保持总质量动量能量和压力正定、怎样避免标签梯度导致虚假自由膨胀，以及再循环后每群内部多束化如何监测。可用三阶矩重建残差、不同标签选择、不同群数及长期磁通演化共同检验；不能只报某个时间的最佳误差。\citep{F635}

递归Hall撕裂也需要闭环验证：同时测每级$L_*,a_*,B_*,\rho_*$和欧姆平衡，在物理耗散固定下细化，并比较静态本征预测与真实片形成史。当前文献已有“次级增长加快”的证据，却没有在同一非线性数据集上验证Hall修正的临界长宽比。强Hall下次级岛缺失仍需区分物理抑制、寿命过短与漏分辨。\citep{L236,L262}

\section{由强导引场新对照明确下来的四项检验}
\paragraph{外层长度由谁选定。}无快速色散波的快重联反例已存在，但匹配率仍含$B_h,w_v$。下一步可固定$d_e,\rho_s$、耗散和分辨率，分别改变流入盒宽、流出盒长、初始片厚、开放／周期边界；逐时测量$w_v$，而非仅拟合峰率。若某种压力或喷流不稳定性选择了长度，应能预测扫描中何时失效。这属于两流体方法本身可推进的问题。\citep{L216,L217}

\paragraph{大岛双重加速能否通过动理学检验。}已有陀螺流体／动理学非线性对照集中在小$\Delta'$，不能代替大$\Delta'$验证。需要采用相同热尺度、边界、种子及足够高次谐波，分别比较第一轮局地电势增强、两轮之间停滞、第二轮排气开口和最终磁通率；再扫描电子回旋半径、Landau响应及微层分辨率。单独重现最后0.1的峰值会漏掉最有区分力的过程。\citep{F672,E475}

\paragraph{热流闭合能否同时预测时间延迟与能量。}非局地六场模型改善温度平滑，却未复现参考动理学中的加热延迟。可以先在冷离子相同几何比较高阶闭合及更完整动理学，再启用热离子六场，测$Q_\parallel$、温度峰值、磁能释放到热响应的时滞和释放磁能归一化的数值误差。有限矩仍不能唯一决定非热尾部；“自由能被闭合移除”也不能直接给出真实粒子熵增。\citep{L244}

\paragraph{三维磁通旁路与实验全局压缩。}局域体力算例显示磁通绕开层，MRX则显示导引场被装置压缩；两者都说明局地片无法独立代表整体系统。先用同样驱动力、相同磁通库存及同样诊断积分域做二维／三维对照，再在装置几何中加入实际边界和驱动。可检验磁通被重联、被绕流带走和暂时堆积的各自比例。Hall／双流体可以研究这些几何效应，电子层加热则仍需更高矩或动理学检验。\citep{L076,F717,L193}
